\chapter{A Typed Solver, Proof Firewall, and TIR Crosswalk}
\label{chap:typed-solver}

\section{Why a solver is needed}

The monograph now contains several different kinds of statements: elementary
identities, classical theorems, exact theorems inside an explicit model,
numerical finite-section evidence, structural TIR assignments, and open bridges.
A prose disclaimer is not enough to guarantee that these categories remain
separate when many relations are composed.  Version 0.7 therefore introduces a
small proof-oriented solver whose primary purpose is not symbolic algebra but
\emph{typed implication control}.

The executable implementation is
\begin{center}
 \texttt{src/secret\_of\_a\_half/identity\_holonomy\_solver.py}.
\end{center}
Its deterministic receipt is generated by
\begin{center}
 \texttt{scripts/run\_identity\_holonomy\_solver.py}.
\end{center}

\section{Claim statuses}

Each solver rule has one of four statuses:
\begin{equation}
 \mathfrak S
 =\{\mathrm{EXACT},\mathrm{STANDARD},\mathrm{MODEL},\mathrm{OPEN}\}.
 \label{eq:solver-status-set}
\end{equation}
Their meanings are:
\begin{description}
 \item[EXACT.] Proved in this repository from the stated definitions or an
 elementary identity.
 \item[STANDARD.] A standard mathematical/physical theorem used with its usual
 assumptions, for example the spinorial double cover or Berry holonomy
 \citep{berry1984,simon1983}.
 \item[MODEL.] A declared TIR/PhaseNav assignment that may be used only in an
 explicitly model-level derivation.
 \item[OPEN.] A missing bridge.  It is never promoted automatically.
\end{description}

This status system is deliberately more restrictive than a generic rule engine.
An OPEN implication remains blocked even when all of its nodes appear elsewhere
in the graph.

\section{Hyperedges rather than one-premise arrows}

Many useful implications require several premises.  For example, exact
cancellation does not follow from $\sigma=1/2$ alone; it also requires a
symmetric detector and a half-turn relative phase.  A rule is therefore a
hyperedge
\begin{equation}
 r=(P,c,s,k,\nu),
 \label{eq:solver-rule}
\end{equation}
where $P=\{p_1,\ldots,p_m\}$ is a finite set of premises, $c$ is the
conclusion, $s\in\mathfrak S$ is status, $k$ is relation kind and $\nu$ is
provenance.  A normalized holonomy coordinate may be attached when the edge is
holonomy-eligible.

Given a fact set $F_0$, exact closure is the least fixed point
\begin{equation}
 F_{n+1}=F_n\cup
 \left\{c_r:
 P_r\subseteq F_n,
 s_r\in\{\mathrm{EXACT},\mathrm{STANDARD}\}
 \right\}.
 \label{eq:exact-closure}
\end{equation}
Model closure additionally permits rules marked MODEL.  OPEN rules are excluded
from both fixed-point operators.

\section{The half-axis rule base}

The v0.7 solver contains rule families corresponding to the principal results
of the monograph:
\begin{align}
 \sigma=1/2
 &\Rightarrow \text{complement fixedness},\\
 \sigma=1/2
 &\Rightarrow H_2=\ln2,\\
 \sigma=1/2+\text{binary-state representation}
 &\Rightarrow \text{Bloch equator},\\
 \text{Bloch equator}+\text{equatorial loop}
 &\Rightarrow \text{Berry holonomy }-1,\\
 \sigma=1/2+\text{symmetric detector}+\text{half-turn}
 &\Rightarrow \text{exact cancellation},\\
 \text{centered zeta chart}+\text{reciprocal chart}
 &\Rightarrow \text{critical-axis set invariance},\\
 \text{projective recurrence}+\text{spin }1/2
 &\Rightarrow \text{spinor double cover}.
 \label{eq:solver-exact-family}
\end{align}
The model family includes
\begin{align}
 \text{binary information}+12\text{ projective cycles}
 &\Rightarrow dI/dq=\ln2/12,\\
 dI/dq=\ln2/12+C=2\pi
 &\Rightarrow \kappa=\ln2/(24\pi),\\
 8\text{ mix sectors}+3\text{ flavours}
 &\Rightarrow 24\text{ declared sectors}.
 \label{eq:solver-model-family}
\end{align}
The first two lines of \cref{eq:solver-model-family} reproduce the crosswalk of
\cref{chap:identity-axis}; their status remains MODEL because the twelve-cycle
assignment is not independently derived here.

\section{The RH firewall}

The most important rule is the one the solver is forbidden to use by default:
\begin{equation}
 \xi(\rho)=0
 +\text{canonical zero-state representation}
 \Longrightarrow
 \text{native closure}.
 \label{eq:solver-open-native}
\end{equation}
This is precisely SOH-C004 from \cref{chap:phasenav-native}.  The consequent
native closure then implies the half-axis by SOH-L009.  The computational
problem is therefore not that the solver cannot find a route to the half.  It
can.  The problem is that one edge on the route has status OPEN.

\begin{criterion}[Solver RH firewall]
\label{crit:rh-firewall}
For any initial fact set that does not contain an independently proved canonical
zero-state theorem, neither exact closure nor explicit model closure may derive
the Riemann Hypothesis.
\end{criterion}

The regression suite checks this criterion.  It also asks the solver for the
missing premises of native closure and requires the result to contain
\texttt{canonical\_zero\_state}.  A future proof can therefore be integrated by
promoting one named dependency rather than rewriting the whole logical graph.

\section{Four numerical root routes}

The solver independently evaluates the four balance conditions of
\cref{thm:four-route-half}.  The complement, entropy-stationarity and
cancellation conditions are solved by bracketed root finding.  The Berry route
checks the $-1$ equatorial holonomy.  The expected receipt is
\begin{equation}
 \sigma_C=\sigma_H=\sigma_I=\sigma_B=0.5
 \label{eq:solver-route-consensus}
\end{equation}
to the declared numerical tolerance.

The point of this computation is not to numerically discover one half.  It is a
cross-implementation test: four functions written from different formulas must
agree with the symbolic theorem and with one another.

\section{Proof traces}

For every derived fact the solver stores the exact hyperedge used, including
premises, status, relation kind, provenance and optional normalized holonomy.
A proof trace is therefore a directed acyclic subgraph of the closure history.
For example, a model trace for $\kappa$ records
\begin{equation}
 \{\text{binary information},12\text{ cycles}\}
 \stackrel{\mathrm{MODEL}}{\Longrightarrow}
 \ln2/12
 \stackrel{C=2\pi}{\Longrightarrow}
 \frac{\ln2}{24\pi}.
 \label{eq:kappa-proof-trace}
\end{equation}
It does not record the result simply as ``derived'' without preserving the
MODEL labels of its premises.

This is useful beyond the present problem.  A long research programme can carry
many correct algebraic implications while still depending on one unproved
physical identification.  Proof traces keep that dependency visible.

\section{TIR cross-references}

The parent Theory of Informational Relations is used here in a sharply bounded
way.  The current TIR/Metatime monograph provides the wider information
normalization, Poincar\'e/phase geometry, flavour architecture and sector
holonomy programme \citep{lipa2026tir}.  The companion PhaseNav repository
provides the geometry-first dependency package used to implement centered
coordinates, normalized cycles, phase-crystal tangents and typed relation edges
\citep{lipa2026phasenav}.

The direction of import is not symmetric:
\begin{enumerate}
 \item Secret of a Half imports TIR assignments only with their existing claim
 status.
 \item Exact half-axis lemmas proved here may be cross-referenced by TIR without
 promoting TIR model assignments to exact physics.
 \item Generated PhaseNav geometry is an executable representation, not proof
 of semantic truth.
 \item The TIR object $W_{ij}$ is not defined or used dynamically here.
\end{enumerate}

\section{Paired implication matrix}

The solver view makes it useful to summarize where implication is genuinely
bidirectional and where it is not.
\begin{longtable}{@{}p{0.28\textwidth}p{0.28\textwidth}p{0.34\textwidth}@{}}
\toprule
Structure A & Structure B & Logical relation\\
\midrule
$\sigma=1/2$ & $\sigma=1-\sigma$ & exact equivalence on $[0,1]$.\\
$\sigma=1/2$ & $H_2=\ln2$ & exact equivalence by strict concavity.\\
$\sigma=1/2$ & Bloch equator & exact equivalence inside binary qubit representation.\\
Bloch equator & Berry $-1$ & implication for the stated equatorial loop; not every loop.\\
$\sigma=1/2$ + half-turn & equal-gain cancellation & exact equivalence under symmetric detector.\\
$\Re s=1/2$ & $\mathcal J(s)=s$ & exact equivalence.\\
centered critical axis & reciprocal image & exact set invariance; not zero-set invariance.\\
projective recurrence & spinor recurrence & standard double-cover ratio $1:2$.\\
$12$ projective cycles & $\ln2$ information cycle & model assignment, not theorem.\\
$8\times3$ & $12\times2$ & exact integer equality; semantic labels model-level.\\
$\xi(\rho)=0$ & native closure & OPEN SOH-C004.\\
native closure & $\Re\rho=1/2$ & exact SOH-L009.\\
\bottomrule
\end{longtable}

The matrix prevents the phrase ``one follows from the other'' from being used
without stating direction and assumptions.

\section{Compatibility with the PhaseNav--Weil programme}

The solver does not replace the Weil, Hermite, prime-tail or adaptive-cutoff
chapters.  Those modules address a different open route: positivity of an
arithmetic operator on an admissible dense family.  The current solver can
represent their eventual promotion in the same typed graph:
\begin{equation}
 \text{complete Weil positivity}
 \stackrel{\mathrm{OPEN}}{\Longrightarrow}
 \text{native closure of zero states}
 \Longrightarrow
 \text{half axis}.
\end{equation}
Until the first arrow is proved, the numerical PSD receipts remain evidence for
finite sections only, exactly as stated in the existing claim ledger.

Likewise, DHSE-001 remains an independent experimental branch of the research
programme.  Version 0.7 does not alter its source, receipts, thresholds or
promotion rules.

\section{Executable receipt}

The solver receipt stores:
\begin{itemize}
 \item the four half-axis route values;
 \item the three exact integer factorizations of $24$;
 \item the conditional numerical value of $\kappa$;
 \item exact closure and its proof trace;
 \item explicit model closure and its proof trace;
 \item the unresolved native-closure bridge;
 \item a Boolean field that must remain false for
 \texttt{riemann\_hypothesis\_derived}.
\end{itemize}
The last field is intentionally a negative test.  A future code change that
turns it true without a corresponding proof-level promotion must fail CI.

\section{Testing strategy}

Version 0.7 adds a focused test module independent of historical deterministic
receipts.  It verifies:
\begin{enumerate}
 \item centered sign and complement about the identity axis;
 \item normalized turns and winding frequency;
 \item spinor sheet parity;
 \item agreement of the four half-axis routes;
 \item positivity of all four off-axis residuals on representative controls;
 \item separation of exact, model and open closure;
 \item exposure of the missing canonical zero-state premise;
 \item consistency of the three $24$ factorizations and the conditional
 $\kappa$ reconstruction.
\end{enumerate}
These tests supplement rather than rewrite the existing 135-test baseline.
Historical receipt failures, if any, are reported separately so a new solver is
not credited or blamed for unrelated archived hashes.

\section{What has actually been gained}

The new solver does not solve RH.  It does something more modest but necessary:
it turns the current architecture into an auditable dependency graph.  The half
is no longer represented merely by a list of suggestive coincidences.  It is a
node selected by several independent exact routes, connected to normalized
recurrence and holonomy, cross-referenced to the wider TIR geometry, and
protected by a machine-readable firewall at the exact point where the proof is
still missing.

This is a useful form of progress because the remaining obstruction is sharply
localized.  A successful future result must either prove the canonical
zero-state/native-closure edge, or supply another exact path that bypasses it
without assuming an equivalent form of RH.

\status{The typed solver and its exact/model firewall are implemented and
reproducible.  The canonical zero-state implication remains OPEN; RH is not
derived.}
